% SPDX-License-Identifier: Apache-2.0
% covers: HostLink, txhdl_parts::remote::host::HostBridge
\datasheet{HostLink}{The wire under a program that drives the design's bus: each ask one Ethernet frame in, each answer one frame out.}

\dsfacts{\code{lib/parts/src/remote/host.rs}}{\code{txhdl\_parts::remote::host::HostLink}, and \code{HostBridge}, it and \code{LiteMaster} joined}%
{\code{//docs:examples}}

\subsection*{Function}

\code{HostLink} is the other direction of \code{RemoteLink}.
There the board's transactions go out to a program that answers them;
here a program on another machine sends an \code{Ask} in a frame, and
\code{HostLink} takes it off the wire and passes it on \code{ask}.
The \code{Answer} that comes back on \code{back} leaves as a frame.
\code{LiteMaster} performs the ask on the design's AXI-Lite bus.
What FireSim's host driver does over PCIe, reading and writing the
simulator's control registers, it can do over this instead (issue
1931).

The frame is \code{RemoteLink}'s, at the same offsets, with two kinds of
its own, so the two directions can share a wire and a program and
neither takes the other's frames for its own.

\code{HostBridge} is \code{HostLink} and \code{LiteMaster} as one unit,
frames on one side and an AXI-Lite host on the other, joined by their
\code{ask} and \code{back} channels inside it, where the lowering
buffers them. Joined by wires in a hand-written top instead, each
channel's valid would depend on its own ready through the other unit, a
combinational loop, which issue 1931 met on the board.

\subsection*{Parameters}

\begin{center}\footnotesize
\begin{tabular}{@{}lp{0.7\textwidth}@{}}
\toprule
Parameter & Meaning \\
\midrule
\code{DEV} & The device number, which every ask this accepts must carry
and every answer it sends carries, as for \code{RemoteLink}. The tables
use 5, as \code{ex\_hostlink} and \code{//firesim/ax7a200b} do. \\
\bottomrule
\end{tabular}
\end{center}

\subsection*{The frame}

27 bytes, before whatever padding the MAC adds to reach the minimum
frame; \code{RemoteLink}'s layout.

\begin{center}\footnotesize
\begin{tabular}{@{}llp{0.52\textwidth}@{}}
\toprule
Bytes & Field & What \\
\midrule
0 to 6 & destination & Broadcast on an answer, since the design does
not know where the program is. \\
6 to 12 & source & \code{02:00:00:00:00:\textit{dev}} on an answer. \\
12 to 14 & type & \code{0x88b5}. \\
14 & kind & \code{0x02} asks, \code{0x82} answers. \\
15 & device & \code{DEV}. \\
16 & tag & The ask's, which the answer repeats. \\
17 & flags & Bit 0 is a write on an ask, an error on an answer. \\
18 to 22 & address & Most significant byte first. \\
22 to 26 & word & The word written, or the word answered. \\
26 & strobe & Which lanes a write means. \\
\bottomrule
\end{tabular}
\end{center}

\dstables{HostLink}

\subsection*{Behaviour}

\begin{itemize}
\item Every byte coming in is taken, so a frame that is not an ask is
read and dropped rather than left to block the MAC.
\item \code{got} counts the bytes of the frame coming in and stops at 63
rather than wrapping, so a long frame's later bytes are never read as
its header or its fields again: frames longer than 64 bytes, such as
ARP, IPv6 or LLDP, are ordinary on a wire.
\item Four bytes are checked as they pass: the two of the type, the
kind and the device. A mismatch sets \code{alien}, which the last byte
clears, so the next frame starts innocent.
\item The tag, the write bit, the address, the word and the strobe are
latched as they arrive, at bytes 16 to 26.
\item An ask goes on \code{ask} at the end of a frame nothing objected
to. A frame that ends by byte 26 is not a whole ask, since its last
field is held a cycle after its byte; the MAC pads every frame to sixty
bytes, so a real ask always goes on past it. With no room on
\code{ask} the ask is dropped, and the program asks again.
\item An answer is taken from \code{back} when no frame is going out,
and its 27 bytes leave a byte per cycle that \code{tx} has room,
spelled from registers by \code{at} rather than kept in a buffer.
\item Taking asks and sending answers are independent.
\end{itemize}

\subsection*{Verification}

\begin{itemize}
\item Eleven tests in \code{lib/parts/src/remote/host.rs}, in
\code{//lib/parts:parts\_test}, run \code{HostLink<5>} and
\code{LiteMaster} with a memory on the bus. Among them:
\code{frames\_that\_are\_not\_this\_devices\_asks\_are\_ignored} sends
another device's ask, another protocol's frame and the other
direction's ask;
\code{a\_long\_ask\_is\_taken\_whole} pads an ask past 64 bytes;
\code{a\_frame\_that\_ends\_at\_the\_strobe\_is\_not\_an\_ask} cuts one
short; \code{a\_long\_frame\_of\_another\_kind\_is\_ignored} hides an
ask's header past byte 64 of another protocol's frame; and
\code{the\_bridge\_is\_the\_two\_units\_joined} sends the same asks
through \code{HostBridge}.
\item \code{//lib/examples:ex\_hostlink} runs a program's asks over a
wire, and the build simulates the three netlists against that run under
nvc and Verilator, among them\newline
\code{//docs:hostlink\_sim\_host\_link\_host\_link\_tb\_test}\newline
\code{//docs:hostlink\_sim\_host\_bridge\_host\_bridge\_tb\_test}\newline
and their \code{vsim} twins.
\end{itemize}

\subsection*{Used by}

\code{ex\_hostlink}, and \code{//firesim/ax7a200b}, where
\code{HostBridge<5>} is the host link to FireSim's GCD design on the
AX7A200B, driven over the Ethernet cable from \code{srv}.

\subsection*{Limits}

One answer frame goes out at a time, and an ask that finds \code{ask}
full is dropped; a program that does not hear an answer asks again.
The frame carries one 32-bit word. Frames carry no authentication, so
any station on the wire that speaks the protocol with this device
number can drive the bus.
